% D06-004. PDR-05, §10.6, §12--13.
\section{Условия дальнейшей адаптации измерительного органа}
\label{sec:ch5_adaptation_limits}
\conceptsource{Статья 5, подраздел 10.6, ограничения и выводы.
Сохранена общая идея изменения любого органа РЗА.
Перспективные способы не названы выполненными опытами.}

\subsection{Выбор изменения и повторная проверка}

Результат сравнения должен определять,
какое ограничение требуется устранить.
Если расхождение связано с токовой чувствительностью,
исследуется зависимость порога от режима.
Если причина относится к поляризации
или сохранению напряжения после коммутации,
изменяется соответствующая часть характеристики.
Аналитическое выражение, адаптивное правило
и обучаемая модель рассматриваются
как возможные реализации одного назначения.

Для нейросетевого органа экспертные примеры
позволяют изменить полученное ранее правило.
Однако увеличение веса отдельной трудной записи
не гарантирует локальности изменения.
Модель может исправить один случай
и ухудшить другие режимы.
Поэтому после коррекции повторяется проверка
как целевого сложного материала,
так и обычных записей, составлявших прежнюю область применения.

Адаптация под энергообъект требует отдельного
контрольного материала, который не участвовал
в выборе изменения, уставки или состояния весов.
Состав измерений, их масштаб, полярность
и доступная предыстория задаются до сопоставления.
Замена входных условий является изменением опыта,
даже если название модели осталось прежним.
Новый результат не переносится автоматически
на прошлые веса или другую схему измерения.

Численный выход сети также не принимается
за вероятность надёжного действия РЗА.
Для такой интерпретации нужна самостоятельная проверка
соответствия оценки наблюдаемым ошибкам.
Иначе высокий выход классификатора
может сопровождать устойчивое неправильное состояние.

\subsection{Уточнение эталона и независимость контроля}

Ручной массив специально содержит сложные
и пограничные режимы.
Его показатели не описывают частоту отказов
на обычном эксплуатационном потоке.
Слепой повтор одним специалистом
проверяет воспроизводимость его решения,
но не заменяет физического разбора спорных случаев.
Из 200 повторов выделены 35 записей
для дополнительного инженерного рассмотрения,
включая четыре случая устойчиво противоположного направления.
Повторная версия метки не заменяет исходную автоматически.

Две точные копии входов между обучением
и валидацией ограничивают независимость
сравнения обученных моделей.
Исключение этих записей из последующего подсчёта
не отменяет их участия в выборе эпохи.
Для окончательной проверки требуется разделение
по группам исходных сигналов
с сохранением истории уже выполненного опыта.

Предусмотренные 54 отложенные осциллограммы
не входят в приведённые валидационные таблицы.
Однако их ответы уже представлены
в отдельно обозначенной части исследовательской галереи.
Если этот материал использовался для выбора модели
или изменения методики, он перестаёт быть
нетронутым итоговым контролем.
Тогда независимая проверка потребует другого набора.

Каждая нейросетевая конфигурация обучена
при одной случайной инициализации.
Повторный выбор записей оценивает
чувствительность фиксированного результата к составу массива,
но не заменяет нескольких независимых обучений.
Различия близких средних оценок
поэтому не определяют универсально лучшую архитектуру.
Для испытания автоматического действия
дополнительно потребуется полная логика БАВР
и оценка времени коммутационного оборудования.

\subsection{Материал для последующего расширения}
\label{sec:ch5_results_pending}
\workinprogress{Физический разбор характерных осциллограмм
и новые варианты адаптивных органов будут добавлены
после авторского уточнения материалов.
Рисунок 11 статьи пока не используется как подтверждение
работы нейросетевого РНМ.
Численные результаты RTDS требуется обновить
после завершения пересмотра эталона.
Спорное распределение групп переключений таблицы 9
остаётся вне окончательных выводов.}

% D06-005 R5-12 BEGIN
\ifdefined\AuthorReview
\par\smallskip\begingroup\small\raggedright\color{red!60!black}
\noindent\textbf{Примечание рецензента [R5-12], содержательное завершение.} Существующие PDR-Q01/Q02 и непринятый рисунок 11 остаются открытыми вопросами самой работы; повторный список вопросов автору не нужен. После их разбора полезнее добавить несколько предметных примеров, чем ещё одну сводную метрику: осциллограмма, причина расхождения, требуемое состояние и следствие для изменения органа. Пока глава убедительно выявляет ограничения и показывает экспертную адаптацию, но не подтверждает завершённое внедрение либо устранение всех выявленных недостатков. Такой промежуточный результат допустим и должен быть назван прямо.\par\endgroup\smallskip
\fi
% D06-005 R5-12 END

\section{Выводы}
\label{sec:ch05_conclusions}
% D06-005 R5-13 BEGIN
\ifdefined\AuthorReview
\par\smallskip\begingroup\small\raggedright\color{red!60!black}
\noindent\textbf{Примечание рецензента [R5-13], редакционная подготовка.} Перед окончательной редакцией оставить подробные методы, но сократить повторные общие объяснения методики в 5.1, 5.6.5, 5.11.1 и выводах. Служебные пометки «Статья 5» скрываются в обычной версии; они не заменяют видимых ссылок на публикации, источники открытых аналитических принципов, архив French/RTE и программу испытаний. Для заимствованных методов поставить предметные ссылки либо точные отсылки к уже имеющемуся обзору. Мелкие подписи общей схемы, матриц и трёхпанельной оценки увеличить по уже существующим заданиям; отдельные сложные задачи на каждую подпись не нужны.\par\endgroup\smallskip
\fi
% D06-005 R5-13 END

\conceptsource{Статья 5, выводы 1--9, в сочетании
с основной постановкой главы. Это проект выводов
по включённому исследовательскому снимку.
Окончательная редакция зависит от физической проверки
эталона и последующих адаптивных органов автора.}

В главе рассмотрен порядок модернизации измерительного
органа РЗА, связывающий физическую постановку,
массовое сопоставление решений, экспертный разбор
и повторную инженерную проверку.
Он применён к органу направления мощности для задач БАВР.
Методика допускает изменение формулы,
адаптивной характеристики либо обучаемого преобразования.
Проверка конкретных РНМ служит предметным примером,
а не утверждением об испытании всех видов защиты.

Массовое сравнение 48\,790 пригодных осциллограмм
показало неоднородность поведения девяти органов
по источникам и токовым режимам.
В группе Open\_EE с характерным током
от 0,01 до 0,05 номинального
средняя пофайловая доля разногласий достигла 52,25\%.
Это выделяет область для проверки чувствительности
и других частей характеристики.
Само разногласие не определяет ошибочный орган,
а согласие большинства не заменяет эталона.

Экспертный материал использован для проверки
формульных решений и адаптации нейросетевых моделей.
Все три компактные конфигурации после экспертного этапа
улучшили распознавание направления относительно
фактически использованных начальных весов.
При этом улучшение направления не всегда сопровождается
улучшением применимости.
Выбор эпохи по той же валидации
и обнаруженные межчастные копии
ограничивают вывод о независимом качестве.

Плотная оценка дополняет выборочные метрики
длительностью ошибок и устойчивостью ответа.
Наименьшее среднее ошибочное время
достигается разными моделями в двух источниках.
Благоприятное среднее может сочетаться
с длительным неправильным состоянием
или с большим числом переключений.
Углубление сети и сокращение спектрального входа
не устранили эти ограничения автоматически.

Дополнительная программа 100 опытов БАВР
позволила рассмотреть зависимость результата
от секции, нагрузки и типа смены направления.
На зафиксированном экспертном эталоне выявлены
пропуски редкой неприменимости,
неустойчивые выходы и несвоевременные устойчивые ответы.
Продолжающийся физический пересмотр меток
ограничивает окончательность этих численных результатов.
Полученные показатели относятся к измерительному органу
и не подменяют испытания всей автоматики.

Основной результат состоит в переходе
от общего процента совпадений
к конкретным условиям недостаточной работы органа.
Они задают направление следующего изменения
и состав повторной проверки.
Такой порядок полезен для совершенствования РЗА
независимо от того, реализована её характеристика
аналитически или средствами машинного обучения.
